\section{Vision y VLM: cambiar la token interface y conservar el interaction pattern}

Lo esencial en la transferencia de Transformer entre dominios no es tratar la imagen «como texto», sino diseñar una nueva tokenización y un nuevo objective. ViT divide la imagen en patches, CLIP alinea los embeddings de imagen y texto, y los VLM conectan después el encoder visual con el modelo de lenguaje mediante un projector, cross-attention o unified tokens.

\subsection{Vision Transformer: un patch es un token visual}

Para una imagen de $H\times W$ y patches de $P\times P$, hay en total
\[
N=\frac{HW}{P^2}
\]
patch tokens. Cada patch se aplana y se proyecta linealmente a $d$ dimensiones; luego se agregan el position embedding y el class token. La complejidad crece con $N^2$, por lo que la alta resolución aumenta rápidamente el attention cost.


\lecturefigure{slide-093.jpg}{ViT methodology: patchify, linear projection y Transformer encoder}{CS25 V5 official deck, p.~93}


\lecturefigure{slide-094.jpg}{ViT vs CNN: sesgo inductivo de la convolución local frente a la interacción global entre tokens}{CS25 V5 official deck, p.~94}

\teachervoice{00:39:11--00:41:34, Karan pasa de la patch tokenization a la comparación entre ViT y CNN, y de ahí a CLIP y los VLM. Lo que el docente enfatiza es cómo un mismo interaction mechanism se extiende mediante distintas input interfaces, no que el preprocesamiento sea igual para todas las modalidades.}

\subsection{CLIP y VLM: el espacio alineado no es el sistema generativo en sí}

CLIP usa una pérdida contrastiva imagen-texto dentro del lote:
\[
\mathcal{L}_{\text{CLIP}}=-\frac{1}{B}\sum_i
\log\frac{\exp(s(v_i,t_i)/\tau)}{\sum_j\exp(s(v_i,t_j)/\tau)}+\text{sym.}
\]
$v_i,t_i$ son los embeddings de imagen y de texto, $s$ es la similitud y $\tau$ es la temperatura. CLIP aprende un espacio de recuperación entre modalidades; un VLM además necesita conectar la representación visual a un decoder generativo y resolver alta resolución, OCR, grounding y hallucination.


\lecturefigure{slide-095.jpg}{CLIP: los encoders de imagen y de texto se alinean en un embedding space compartido}{CS25 V5 official deck, p.~95}


\lecturefigure{slide-097.jpg}{VLM basics: vision encoder, adapter/projector y language model}{CS25 V5 official deck, p.~97}

\begin{knowledgebox}{Términos clave: patch, encoder, projector y grounding}
Un \textbf{patch token} es el vector de un bloque local de la imagen; el \textbf{vision encoder} extrae la representación visual; el \textbf{projector/adapter} proyecta la dimensión visual a la interfaz del modelo de lenguaje; la \textbf{cross-attention} permite que un conjunto de tokens consulte a otro; el \textbf{grounding} exige que el texto se alinee con entidades o regiones de la imagen. La slide 97 dibuja estos módulos como un pipeline de VLM. En la lectura de la figura conviene revisar el número de tokens visuales, la resolución, qué partes se congelan y cuáles se entrenan, y la loss de lenguaje; que los módulos estén conectados no significa que el modelo realmente localice la evidencia.
\end{knowledgebox}

\begin{knowledgebox}{Relación jerárquica entre ViT, CLIP y VLM}
ViT es una arquitectura de encoder visual; CLIP es un paradigma de entrenamiento contrastivo imagen-texto; un VLM es un sistema capaz de comprender o generar de forma conjunta contenido visual y lingüístico. Un VLM puede usar ViT y un preentrenamiento tipo CLIP, pero los tres términos no son sinónimos.
\end{knowledgebox}

\subsection{Resumen de la sección}

La extensión entre modalidades cambia primero la token interface y, en segundo lugar, el objective de preentrenamiento y la forma de alineación. ViT, CLIP y los VLM resuelven, respectivamente, la codificación visual, la representación entre modalidades y la interacción generativa; el costo de la alta resolución, el grounding y la hallucination todavía requieren una evaluación independiente.
